\documentclass[11pt,reqno]{amsart}
\usepackage[T1]{fontenc}
\usepackage[utf8]{inputenc}
\usepackage{lmodern}
\usepackage{amsmath,amssymb,amsfonts,amsthm,mathtools,mathrsfs,bm}
\usepackage{enumitem}
\usepackage{booktabs}
\usepackage{microtype}
\usepackage[colorlinks=true,linkcolor=blue,citecolor=blue,urlcolor=blue]{hyperref}
\usepackage[nameinlink,capitalize]{cleveref}
\allowdisplaybreaks
\setlength{\emergencystretch}{2em}

\newtheorem{theorem}{Theorem}[section]
\newtheorem{proposition}[theorem]{Proposition}
\newtheorem{lemma}[theorem]{Lemma}
\newtheorem{corollary}[theorem]{Corollary}
\newtheorem{definition}[theorem]{Definition}
\newtheorem{assumption}[theorem]{Assumption}
\newtheorem{example}[theorem]{Example}
\theoremstyle{remark}
\newtheorem{remark}[theorem]{Remark}

\newcommand{\E}{\mathbb E}
\newcommand{\Pp}{\mathbb P}
\newcommand{\Qp}{\mathbb Q}
\newcommand{\R}{\mathbb R}
\newcommand{\N}{\mathbb N}
\newcommand{\T}{\mathbb T}
\newcommand{\cE}{\mathcal E}
\newcommand{\cP}{\mathcal P}
\newcommand{\cL}{\mathcal L}
\newcommand{\Id}{\mathrm{Id}}
\newcommand{\tr}{\operatorname{tr}}
\newcommand{\Var}{\operatorname{Var}}
\newcommand{\Cov}{\operatorname{Cov}}
\newcommand{\osc}{\operatorname{osc}}
\newcommand{\TV}{\mathrm{TV}}
\newcommand{\norm}[1]{\left\lVert #1\right\rVert}
\newcommand{\abs}[1]{\left\lvert #1\right\rvert}
\newcommand{\pair}[2]{\left\langle #1,#2\right\rangle}

\title[Physical-time deterministic homogenization]{Physical-Time Homogenization in Deterministic Driven Collision Ensembles}
\author{Qian Qi}
\address{Peking University, Beijing 100871, China}
\email{qiqian@pku.edu.cn}
\subjclass[2020]{Primary 60F17, 37A50; Secondary 60H10, 37C30}
\keywords{deterministic homogenization, renewal time change, rough paths, driven ensemble, Green--Kubo covariance}
\date{August 29, 2026}

\begin{document}
\raggedbottom
\begin{abstract}
We derive a stochastic physical-time limit from the deterministic driven maps
of Paper I.  The only probability is Lebesgue uncertainty of the initial
microstate after the macro current preparation.  Four branch impulses are
chosen as actual collision-port momentum transfers and are centered for every
driven current.  Their collision covariance and the nonconstant branch roof
are explicit:
\[
 C(a)=\operatorname{diag}(1-a,6(1+a)),\qquad
 \bar\tau(a)=\frac{3-a}{2}.
\]
The canonical branch sequence is iid because the deterministic full-cover map
resets the future coordinate.  We prove the enhanced invariance principle,
renewal time change, physical covariance, coefficient response, and an
exponential random-horizon estimate strong enough for entropic dynamic
programming.  No Brownian noise, Markov transition matrix, or stochastic clock
is inserted at the microscopic level.
\end{abstract}
\maketitle
\tableofcontents

\section{The driven deterministic platform}
Fix a macro current $a\in(-1,1)$ and let $B_a,m_a,w(a)$ be as in Paper I.
Put
\[
 d_1=(2,0),\quad d_2=(-1,0),\quad
 d_3=(0,4),\quad d_4=(0,-3),
\]
and choose physical flight times
\[
 \tau_1=\tau_2=2,\qquad \tau_3=\tau_4=1.
\]
The vectors are branchwise collision-port impulses and the roof is the
branchwise duration of the deterministic suspension over $B_a$.

\begin{proposition}[Exact centering and characteristics]\label{prop:chars}
For every $a\in(-1,1)$,
\[
 \sum_iw_i(a)d_i=0.
\]
More strongly, with $\tau_i\in\{1,2\}$,
\[
 \E_{m_a}[d_{I_0}\mid\tau_{I_0}]=0.
\]
Furthermore,
\[
 C(a)=\sum_iw_i(a)d_i\otimes d_i
 =\begin{pmatrix}1-a&0\\0&6(1+a)\end{pmatrix},
\]
and
\[
 \bar\tau(a)=\sum_iw_i(a)\tau_i=\frac{3-a}{2}.
\]
On every compact $K\Subset(-1,1)$, $C(a)$ is uniformly positive definite.
\end{proposition}
\begin{proof}
Within the negative-current branches the weight ratio is $1:2$, so
$1(2,0)+2(-1,0)=0$.  Within the positive-current branches the ratio is $3:4$,
so $3(0,4)+4(0,-3)=0$.  These are exactly the two roof classes, proving the
conditional centering.  Direct summation gives the covariance and roof.
\end{proof}

\section{Collision-clock rough limit}
Let $I_k$ be the branch process of $B_a$ under $m_a$ and define
\[
 W_n(t)=n^{-1/2}\sum_{k<nt}d_{I_k}.
\]
Let $\mathbb W_n$ be its canonical polygonal second level.

\begin{theorem}[Deterministic enhanced invariance principle]\label{thm:rough}
For every $p>2$,
\[
 (W_n,\mathbb W_n)\Longrightarrow(B,\mathbb B^{\rm Str})
\]
in $p$-variation, where $B$ is Brownian motion with covariance $C(a)$.  The
antisymmetric area anomaly is zero.
\end{theorem}
\begin{proof}
By the cylinder formula of Paper I, the branch symbols are iid with law $w(a)$.
The increments are bounded and centered by \cref{prop:chars}.  The martingale
functional CLT and discrete Burkholder estimates give first- and second-level
tightness; the polygonal diagonal identifies the Stratonovich lift
\cite{HallHeyde1980,FrizHairer2014}.  Independence makes every off-diagonal
predictable antisymmetric compensator vanish.
\end{proof}

\section{The actual physical clock}
Let
\[
 S_n^\tau=\sum_{k=0}^{n-1}\tau_{I_k},\qquad
 N_a(t)=\max\{n:S_n^\tau\le t\}.
\]

\begin{lemma}[Renewal concentration]\label{lem:renewal}
For every $r\ge2$ and $T<\infty$,
\[
 \E\max_{m\le n}\abs{S_m^\tau-m\bar\tau(a)}^r
 \le C_rn^{r/2},
\]
and there are $c,C>0$ such that
\[
 \Pp\left(
 \abs{N_a(t)-t/\bar\tau(a)}>u\sqrt t
 \right)\le Ce^{-cu^2}
\]
for $t\le T\varepsilon^{-2}$ and admissible $u$.
\end{lemma}
\begin{proof}
The roof variables are iid, bounded in $[1,2]$, and centered after subtracting
$\bar\tau$.  Burkholder or Hoeffding maximal estimates give the first bound.
Monotone inversion and bounded increments give the sub-Gaussian inverse-clock
bound.
\end{proof}

\begin{theorem}[Physical-time rough limit]\label{thm:physical}
With physical scaling
\[
 \widetilde W_n(t)=n^{-1/2}\sum_{k<N_a(nt)}d_{I_k},
\]
one has
\[
 (\widetilde W_n,\widetilde{\mathbb W}_n)
 \Longrightarrow
 \left(\bar\tau(a)^{-1/2}B,
       \bar\tau(a)^{-1}\mathbb B^{\rm Str}\right).
\]
Thus
\[
 A_{\rm phys}(a)=\frac{C(a)}{\bar\tau(a)}
 =\begin{pmatrix}
 \dfrac{2(1-a)}{3-a}&0\\[2mm]
 0&\dfrac{12(1+a)}{3-a}
 \end{pmatrix}.
\]
\end{theorem}
\begin{proof}
The uniform law of large numbers for $N_a(nt)/n$ and the collision rough limit
combine through continuity of monotone rough time change.  Since the impulse
mean is zero, there is no clock-drift correction.
\end{proof}

\section{A deterministic slow system}
Let $b:\R^2\to\R^2$ be globally Lipschitz.  At collision $k$ define
\[
 X_{k+1}^\varepsilon
 =X_k^\varepsilon+\varepsilon d_{I_k}
  +\varepsilon^2\tau_{I_k}b(X_k^\varepsilon),
\qquad
 t_{k+1}^\varepsilon=t_k^\varepsilon+
 \varepsilon^2\tau_{I_k}.
\]

\begin{theorem}[Physical diffusion from deterministic collisions]\label{thm:sde}
The physical-time interpolation converges weakly to
\[
 dX_t=b(X_t)\,dt+\sigma(a)\,dW_t,
 \qquad \sigma(a)\sigma(a)^T=A_{\rm phys}(a).
\]
The convergence holds jointly with the canonical rough lift.
\end{theorem}
\begin{proof}
Collision-clock convergence follows from \cref{thm:rough} and continuity of
the It\^o--Lyons map.  The drift Riemann sum contains
$\varepsilon^2\tau_{I_k}b$, so division by physical time yields exactly $b$.
Apply \cref{thm:physical}.
\end{proof}

\section{Coefficient response}
\begin{corollary}[Endogenous-ensemble response]\label{cor:response}
The coefficients are analytic in $a\in(-1,1)$.  In particular,
\[
 \partial_a C(a)=\operatorname{diag}(-1,6),
 \qquad \bar\tau'(a)=-\frac12,
\]
and $A_{\rm phys}'(a)$ is nonzero everywhere.
Through Paper I, $a$ is equivalent to the derived conjugate parameter
$\theta(a)=I'(a)$, with $d\theta/da=(1-a^2)^{-1}$.
\end{corollary}

\section{Exponential random-window control}
Let
\[
 n_\varepsilon(T)=\left\lfloor
 \frac{T}{\varepsilon^2\bar\tau(a)}\right\rfloor,
 \qquad N_\varepsilon(T)=N_a(T/\varepsilon^2).
\]

\begin{theorem}[Exponential random-horizon estimate]\label{thm:exp-window}
For every fixed $\lambda,T>0$ there is $C_{\lambda,T}$ such that
\[
 \log\E\exp\left
 \{\lambda\abs{
 \varepsilon\sum_{k=n_\varepsilon(T)}^{N_\varepsilon(T)-1}d_{I_k}}
 \right\}
 \le C_{\lambda,T}\sqrt\varepsilon.
\]
Moreover,
\[
 \log\E\exp\left
 \{\lambda\varepsilon^2
 \abs{N_\varepsilon(T)-n_\varepsilon(T)}\right\}
 \le C_{\lambda,T}\varepsilon.
\]
\end{theorem}
\begin{proof}
Condition first on the complete roof/sign sequence.  By
\cref{prop:chars}, the subtype impulses remain independent and centered inside
each roof class.  Hence for every deterministic window of length $m$,
\[
 \E\left[
 e^{\xi\cdot\sum d_{I_k}}\mid(\tau_{I_k})
 \right]\le e^{C\abs\xi^2m}.
\]
The inverse-clock deviation has the sub-Gaussian tail of
\cref{lem:renewal}.  Split into dyadic window lengths and sum; typical length
is $O(\varepsilon^{-1})$, so the scaled displacement has exponential Orlicz
size $O(\sqrt\varepsilon)$.  The second estimate follows directly from the
inverse-clock mgf.
\end{proof}

\section{Main theorem}
\begin{theorem}[Deterministic-to-stochastic physical closure]\label{thm:main}
For every macro current $a\in(-1,1)$, the same deterministic driven collision
map supplies the natural path law, centered mechanical impulses, nonconstant
physical roof, Brownian rough limit, physical diffusion, coefficient response,
and exponential random-horizon control.  No microscopic stochastic forcing is
used.
\end{theorem}
\begin{proof}
Combine \cref{prop:chars,thm:rough,thm:physical,thm:sde,cor:response,thm:exp-window}.
\end{proof}

\section*{Scope and review status}
The theorem is a complete internal proof in the declared deterministic
full-cover collision scope.  It does not claim that an arbitrary periodic
Sinai billiard has the same exact symbolic reset or that independent external
review has occurred.

\bibliographystyle{plain}
\bibliography{references}
\end{document}
